\section{分层评价设计：检验接口与证据来源约束}

\begin{figure}[H]
  \centering
  \includegraphics[width=0.98\columnwidth]{fig06_evaluation_design}
\caption{分层评价设计与检索协议。在固定的资料库、查询集和前 $k$ 名规则下，评价分别核查测量侧观测包的可回溯性，以及资料扩展、查询、语言、编码器和分块方式对前 10 个结果中核心资料占比的影响。}
  \label{fig:evaluation-design}
\end{figure}

本章评价测量—证据接口的可追溯性及其证据侧检索行为，不报告自动故障诊断、根因识别或生成文本的工程正确性。评价对象分为两个层次：测量侧检查一份观测包能否回溯到确定的输入、参数和派生产物；证据侧检查在同领域资料逐步加入时，仅依赖语义相似度的检索能否仍让预先指定的核心资料保留在前 10 个结果中，并保持足够占比（下文简称“权威资料可见性”）。候选解释的工程适用性需要运行/维护上下文和领域专家复核，不作为本章的自动化评价对象。

\FloatBarrier
\subsection{评价对象与检索协议}

观测层核查观测包的可回溯性：检查活跃记录、时间排序、分析参数、工具输出与记事板字段能否相互对应；对代表性追溯链，再核对完整报告和紧凑告警字段中的关键描述能否回指到脚本输出及时间—通道位置。第五章据此展示一条完整链条。

证据侧使用第 3 章定义的 64 份技术资料：13 份 IAEA 资料构成预设核心集合 $\mathcal{D}_{\mathrm{core}}$，51 份核仪表与控制同领域资料构成补充集合 $\mathcal{D}_{\mathrm{sup}}$。补充资料扩展技术覆盖，核心资料对应预设权威来源；来源要求是资料范围扩展后，前 10 个结果仍应保留足够比例的核心资料。

令 $\mathcal{D}_{\mathrm{sup}}^{(m,t)}$ 表示第 $t$ 次试验中抽取的 $m$ 份补充资料，活跃资料池定义为

\[
\mathcal{D}_{m,t}=\mathcal{D}_{\mathrm{core}}\cup\mathcal{D}_{\mathrm{sup}}^{(m,t)},
\qquad m\in\{0,5,10,\ldots,51\}.
\]

当 $m=0$ 时，仅检索核心资料；当 $m>0$ 时，对同一规模的补充资料进行固定随机种子下的重复抽样。每次检索返回前 10 个段落。该设计不把“资料扩展”简单等同于添加噪声，而是观察相关补充资料加入后，排名与权威资料可见性如何共同变化。

查询集由人工固定和复核。\texttt{nomic-embed-text} 字符分块扩库协议使用中英文核心术语与测量短语；MiniLM encoder 对照另加入由 observation packet 保留线索形成的技术细节表述，如重复短时偏离、同步零值段和高通道共变而无明显滞后，用于检验结构化观察的表述方式是否改变权威资料可见性。对于 \texttt{nomic-embed-text} 协议，字符分块扩库实验覆盖 17 个有效的长度—重叠率配置：长度为 800、1000、1200、1500 字符，重叠为 50、80、100、120、150 字符；每个非零扩库规模下进行 100 次补充资料抽样。该协议考察资料范围、语言和分块边界，并复用其索引进行 self-retrieval。

MiniLM 对照在相同资料库、top-10 规则和来源指标下由三项互补协议组成。第一项仅使用多语种 \nolinkurl{sentence-transformers/paraphrase-multilingual-MiniLM-L12-v2}（MiniLM-L12）：为避免将中文 query 与英文向 encoder 的语言失配混入 query 效果，在 800/80、1200/120 和 1500/150 三种字符分块下比较简短 query 与由 observation packet 线索形成的技术细节 query（图~\ref{fig:query-effect}）。第二项将 MiniLM-L12 与通用 \nolinkurl{sentence-transformers/all-MiniLM-L6-v2}（MiniLM-L6）置于相同中英文 query 集和三种字符分块设置下，考察 query、语言、encoder 与分块的交互（图~\ref{fig:configuration-interaction}）。两模型的训练目标、训练数据和语言覆盖均不同，故该对照比较实际 encoder 配置，而非单独检验 Transformer 层数。第三项采用 240 词、24 词重叠的固定词数分块，用于核查权威资料可见性下降的风险是否只出现在 \texttt{nomic-embed-text} 的字符分块中。各 MiniLM 非零扩库条件进行 10 次补充资料抽样；不同 encoder 或分块单位的余弦分数不直接横向比较，比较重点是各协议内权威资料可见性随资料扩展的变化。

\subsection{指标、检索单元与结果解释范围}

本研究以归一化向量的余弦相似度作为无显式证据来源约束的 dense-retrieval baseline；其最高分仅反映语义近邻行为。对于检索列表 $R_k(q)=\{d_1,\ldots,d_k\}$，定义权威资料可见性为

\[
\operatorname{CorePriority@}k(q)=
\frac{1}{k}\sum_{i=1}^{k}
\mathbb{I}\!\left(\operatorname{doc}(d_i)\in\mathcal{D}_{\mathrm{core}}\right).
\]

当 $k=10$ 时，值为 0.6 表示前 10 个结果中有 6 个来自核心资料，值为 0 则没有核心资料。该指标只衡量检索结果对预设来源要求的符合程度，不代表段落相关性、语义相似度或工程正确性；同时记录第一条核心来源段落的排名。

为核查分块是否仍能找回已知段落，对 IAEA 核心资料进行自检索：以段落首句、标题、随机窗口或领域术语构造查询，分别报告严格同分块和同文档匹配下的前十名召回率与平均倒数排名。前者要求找回同一段落，后者允许返回同文档相邻分块；该核查用于解释检索单元与分块边界，不等同于权威资料可见性评价。

本章结果限定于固定资料库、预先固定的 query 协议和既定检索流程；检索过程不使用语言模型即时改写 query。结果刻画资料范围、查询表述、语言、encoder 和分块方式如何改变前 10 个结果中的核心资料占比，以及该变化是否与最高余弦相似度背离；不同语言、encoder 或分块方式不据此作普适排名。第五章据此呈现观测可追溯链与权威资料可见性结果。
